
\documentclass[11pt,a4paper]{article}

% Page layout
\usepackage[a4paper, margin=1.5cm]{geometry}

% Font and encoding
\usepackage[T1]{fontenc}
\usepackage[utf8]{inputenc}
\usepackage{lmodern}

% Math packages
\usepackage{amsmath}    % Basic math symbols
\usepackage{amssymb}    % Additional math symbols
\usepackage{amsfonts}   % Math fonts
\usepackage{amsthm}     % Theorem environments
\usepackage{mathtools}  % Extensions to amsmath
\usepackage{bm}         % Bold math symbols
\usepackage{cancel}     % Cancel out terms in equations

% Graphics and figures
\usepackage{graphicx}   % Include graphics

\usepackage[font=small,labelfont=bf,labelsep=quad,format=line]{caption}
%\captionsetup{font=small,format=line}

\usepackage{float}      % Control float placements
\usepackage{wrapfig}    % Wrapping text around figures
\usepackage{subfig}
\usepackage{calc}

% Tables
\usepackage{booktabs}   % Professional-looking tables
\usepackage{array}      % Custom table column formatting
\usepackage{multirow}   % Multi-row cells in tables
\usepackage{longtable}  % Tables across pages
\usepackage{tabularx}   % Tables with fixed-width columns

% Hyperlinks

% Bibliography (for BibTeX or BibLaTeX use)
% Uncomment below to use BibLaTeX, and set backend=biber for UTF-8 compatibility

\usepackage[backend=biber,style=numeric,sorting=none]{biblatex}
\addbibresource{report1.bib}
\bibliographystyle{plain}
\bibliography{report1}

% Appendix
\usepackage[toc,page]{appendix}

% Formatting
\usepackage{setspace}
%\onehalfspacing         % Set line spacing to 1.5
\usepackage{titlesec}   % Customize section titles
\titleformat{\section}{\normalfont\Large\bfseries}{\thesection}{1em}{}
\titleformat{\subsection}{\normalfont\large\bfseries}{\thesubsection}{1em}{}
\titleformat{\subsubsection}{\normalfont\normalsize\bfseries}{\thesubsubsection}{1em}{}

\sectionfont{\normalsize} % Set section titles to large size
\subsectionfont{\normalsize} % Set subsection titles to normal size
\subsubsectionfont{\small} % Set subsubsection titles to small size

\usepackage[ruled]{algorithm2e}

% References

\usepackage{hyperref}
\usepackage{cleveref}   % Smart referencing with \cref and \Cref
\usepackage{sectsty}



% Title page
\title{Signature of Variable Binding in Episodic Memory}
\author{Keith Adi Grace and Tim}

\date{\today}


\begin{document}

\maketitle
\abstract{Brain binding}

\section{Introduction}%
  \label{sec:Introduction}

\subsection{Reconciling Episodic Memory with Working Memory Solutions}%
  \label{sub:Reconciling Episodic Memory with Working Memory Solutions}

 Considering the temporal structure of the task, an alternative hypothesis for task solution is through working memory. Specifically, the task can be reinterpreted as a cued working memory retrieval task, in which presentation of the example sequence serves as a cue for retrieval of the first or the second item in the test sequence. This setting can be mapped directly to a working memory control task where macaques are trained to reproduce one of two simultaneously presented items after a delay, depending on a cue that either precedes or follows stimulus presentation (coined "pro" and "retro" conditions) \cite{panichelloSmucwma2021}. For the retro task, macaques learn a solution in which the two stimuli are represented independently in orthogonal PFC neural activity subspaces before selection, and get aligned based on task relevance after cue presentation. In pro trials, alignment is immediate and the subspaces used to represent the selected/unselected items are the same ones used in the retro task. Importantly, the authors establish a correspondence between the transformation of working memory contents used in the two tasks with cross-decoder analyses\cite{panichelloSmucwma2021}. \\

 Generally, this idea that items are stored in orthogonal subspaces in the working memory and context-dependent transformations are used to route relevant representations to output-potent spaces is a common theme in the working memory literature \cite{xieGswmmpc2022,libbyRdrismr2021,liFULRSWMMPC2025,panichelloSmucwma2021}. A recent theoretical proposal \cite{whittingtonttaSsepsmauhcm2025}.

\section{Results}%
  \label{sec:Results}

  \subsection{Representation Similarity Analysis}%
    \label{sub:Representation Similarity Analysis}
    

\section{Methods}%
  \label{sec:Methods}

  \subsection{RSA Analysis}%
    \label{sub:RSA Analysis}
   
  
  
  
% \input{em_wm_motivation.tex}
\printbibliography
\setcounter{figure}{0}
\renewcommand{\thefigure}{S\arabic{figure}}
\renewcommand{\figurename}{Supplementary Figure}

\setcounter{subsection}{0}
\renewcommand{\thesubsection}{S\arabic{subsection}}
\renewcommand{\thesubsubsection}{S\arabic{subsection}.\arabic{subsubsection}}

% new page
\clearpage

%\input{appendix.tex}


\end{document}
